\chapter{绪论}
\label{cha:intro}

\section{引言}

%  进ai前的原始提示词

% gnss/ins惯性导航
% 近些年 很多应用场景对定位服务有着要求 涉及到智能驾驶 船舶航行 导航系统 
% 以及军事 测绘等相关应用场景，都有着相当广泛的应用。传统实现这些需求的方法一般是借助
% GNSS（英文解释）全球定位系统来进行帮助，然而其精度受到多种因素限制，一种是GNSS使用环境
% 带来的天然误差，信号穿过对流层和电离层会有必然的折射、在城市的遮挡环境下，建筑会形成多径效应甚至
% NLOS（英文解释）， 即便电磁波可以穿过障碍物（如树木），也会受到衰减（这里的引用待补充）。
% 此外，在GNSS的诸多干扰也和工程实践中的场景有着联系，例如GNSS用户端和卫星端授时误差的不确定性
% 用户端接受设备的快速移动带来的不可忽略的多普勒（英文释义）效应，都会对GNSS的定位精度造成一定影响。

% 一种结合INS（英文释义）惯性传感器的捷联惯导系统（补充一些），借助惯性元件的（什么什么）特性
% ，补齐了上述GNSS定位的某些短板，当然INS器件也有着自己的误差来源，且公式相对来说计算复杂，这主要体现在要使用INS
% 必须考虑到地球自转，重力方向，海拔因素对地理/导航坐标系的影响，且INS由于传感器本身的性质（这里的引用待补充），
% ，加速度计的零偏，加速度计标定刻度，由安装方向带来的失准角和不严格正交的特性，都可能会带来额外的误差。另外，在公式计算
% 中，在计算地球形状、计算小角度旋转的反对称矩阵时，可能会存在一些数学模型的近似，在某些情况下也会带来误差。

% 综上所述，对于城市空间驾驶、定位的一般平面上导航问题来说，如何运用合适的算法融合INS数据和GNSS数据，成为解决问题的关键


\section{引言}

近年来，随着信息技术与智能系统的快速发展，越来越多的应用场景对高精度定位与导航服务提出了更高要求。在智能驾驶领域，车辆需要实时获取自身位置、速度及姿态信息，以实现车道级导航、自动驾驶控制以及复杂交通环境下的安全决策；在船舶航行与航空导航领域，精确的位置信息对于航迹规划、避碰控制以及长时间稳定运行具有重要意义；在测绘与地理信息系统领域，高精度定位技术被广泛应用于地形测量、工程建设以及灾害监测等任务；此外，在军事侦察、无人系统导航以及应急救援等场景中，稳定可靠的导航能力同样具有重要意义。因此，研究高可靠、高精度的定位与导航方法，已经成为导航领域的重要研究方向之一。
目前，实现上述定位需求最常用的方法之一是依赖全球导航卫星系统（Global Navigation Satellite System，GNSS）。GNSS 能够通过接收来自多颗导航卫星的信号，实现全球范围内的三维定位与测速功能，在开阔环境中通常能够提供较为稳定且精度较高的位置与速度信息。然而，在实际应用中，GNSS 的定位精度仍然受到多种因素的影响。首先，GNSS 信号在传播过程中需要穿过对流层与电离层，大气折射会引入额外的传播延迟误差，从而影响测距精度。此外，在城市峡谷、隧道出入口以及高层建筑密集区域等复杂环境中，建筑物的遮挡会使卫星信号发生反射，从而产生多径效应（Multipath Effect）甚至非视距传播（Non-Line-of-Sight，NLOS）现象，这类误差通常具有较强的不确定性，难以通过简单模型进行补偿。与此同时，即便电磁波能够穿过部分障碍物（如树木或轻质结构），信号强度仍会发生衰减，从而降低接收质量并影响定位稳定性。
除环境因素外，GNSS 系统本身也存在一定误差来源。例如，卫星端与用户接收端之间的授时误差会直接影响测距结果；在用户设备高速运动的情况下，还会产生显著的多普勒效应（Doppler Effect），对载波信号频率产生影响，从而对速度与位置解算带来额外误差。在极端情况下，当卫星信号受到严重遮挡或干扰时，GNSS 甚至可能出现短时间中断，导致导航系统无法持续输出可靠的定位信息。
为弥补 GNSS 在复杂环境下易受干扰、可用性不足的缺点，惯性导航系统（Inertial Navigation System，INS）逐渐成为 GNSS 的重要补充。捷联式惯性导航系统（Strapdown Inertial Navigation System，SINS）利用安装在载体上的加速度计与陀螺仪等惯性传感器，能够在不依赖外部信号的情况下，通过对角速度与比力进行积分运算，实时估计载体的速度、位置及姿态信息。相比 GNSS，INS 具有输出频率高、短时精度高以及抗干扰能力强等优点，因此在 GNSS 信号短时间失锁或质量下降时，仍能够维持系统的连续导航能力，在组合导航系统中发挥着关键作用。
然而，INS 也存在其固有缺陷。由于惯性导航系统的计算过程本质上依赖积分运算，传感器的微小误差会随着时间逐渐累积，从而导致导航误差呈发散趋势。在数学建模过程中，需要同时考虑地球自转、地球曲率、重力方向以及海拔高度变化等因素对导航坐标系的影响，使得系统模型相对复杂。此外，由于惯性传感器自身的物理特性，加速度计与陀螺仪通常存在零偏误差、比例因子误差以及随机噪声等问题；在实际安装过程中，还可能产生失准角误差以及非严格正交误差，这些误差均会在积分过程中不断放大，从而影响导航精度。同时，在算法实现过程中，对地球形状的建模以及小角度旋转反对称矩阵的近似处理，也可能在特定条件下引入额外的建模误差。
综上所述，在典型的城市道路导航或地面运动载体定位问题中，单一导航系统往往难以同时满足长期稳定性与短期高精度的需求。GNSS 具有长期误差不发散的优点，但容易受到环境干扰；INS 则具有短时精度高和连续性强的特点，但误差会随时间迅速累积。因此，如何利用合适的数据融合算法，将 GNSS 的长期稳定性与 INS 的短期高频特性进行有效结合，成为当前组合导航系统研究中的关键问题之一。

\subsection{惯性传感器的数学模型}

\subsubsection{直接建模方法}

在基于滤波的GNSS/INS组合导航系统中，根据系统状态变量选取方式的不同，
常见的建模方法可分为直接方法（Direct Method）和间接方法（Indirect Method）。
其中，直接方法是一种以导航参数本身为系统状态量的建模方式，
该方法在早期组合导航研究中得到了广泛应用。

直接建模方法的基本思想是：将惯性导航系统计算得到的导航参数，
例如位置、速度和姿态等作为滤波器的状态变量，
并利用惯性导航的运动学方程描述系统状态的演化过程。
同时，GNSS观测量（如位置和速度）直接作为观测方程，
用于修正INS计算结果，从而实现组合导航解算。

与间接方法相比，直接方法具有较为直观的物理意义。
该方法将导航参数本身作为状态变量，
使滤波器不仅承担误差估计的功能，
同时也直接完成导航解算过程，
从而避免了在惯性导航计算与误差修正之间进行重复计算，
在一定程度上提高了计算效率。
此外，由于直接方法采用完整的导航运动学模型，
其状态方程能够较准确地反映系统真实的动态特性，
相比间接方法中对误差方程进行一阶线性近似的处理，
具有更高的物理一致性。 \cite{HU2018755}

在直接建模方法中，系统状态通常定义为导航参数向量，例如

\begin{equation}
\mathbf{x}
=
\begin{bmatrix}
\boldsymbol{p} \\
\boldsymbol{v} \\
\boldsymbol{\phi}
\end{bmatrix}
\end{equation}

其中，
$\boldsymbol{p}$ 表示位置向量，
$\boldsymbol{v}$ 表示速度向量，
$\boldsymbol{\phi}$ 表示姿态角向量（如欧拉角）。

系统状态方程通常由惯性导航的姿态方程、
速度方程和位置方程构成。
例如，在地理坐标系（East-North-Up, ENU）下，
姿态更新方程可表示为

\begin{equation}
\dot{\boldsymbol{\phi}}
=
\mathbf{M}
\left(
\boldsymbol{\omega}_{gb}^{b}
+
\boldsymbol{\varepsilon}
\right)
+
\mathbf{M}
\boldsymbol{w}_{\varepsilon}
\end{equation}

其中，
$\boldsymbol{\omega}_{gb}^{b}$ 表示机体系到导航坐标系的角速度，
$\boldsymbol{\varepsilon}$ 表示陀螺零偏，
$\boldsymbol{w}_{\varepsilon}$ 表示陀螺白噪声，
$\mathbf{M}$ 为姿态变换矩阵。 \cite{HU2018755}

对应的速度方程可表示为

\begin{equation}
\dot{\boldsymbol{v}}
=
\boldsymbol{f}
-
\left(
2\boldsymbol{\omega}_{ie}
+
\boldsymbol{\omega}_{en}
\right)
\times
\boldsymbol{v}
+
\boldsymbol{g}
\end{equation}

位置方程则可表示为

\begin{equation}
\dot{\boldsymbol{p}}
=
\mathbf{T}
\boldsymbol{v}
\end{equation}

上述方程共同构成了直接建模方法中的系统状态模型，
其通过惯性测量单元（IMU）输出的角速度和比力信息，
递推计算导航参数。

% 在观测模型方面，
% GNSS提供的位置和速度信息通常作为观测量，
% 其观测方程可表示为

% \begin{equation}
% \mathbf{z}
% =
% \mathbf{H}
% \mathbf{x}
% +
% \mathbf{v}
% \end{equation}

% 其中，
% $\mathbf{z}$ 表示GNSS观测向量，
% $\mathbf{H}$ 为观测矩阵，
% $\mathbf{v}$ 表示观测噪声。

需要指出的是，
由于直接方法的状态方程通常包含非线性项，
例如姿态更新和速度更新中的非线性耦合关系，
传统线性卡尔曼滤波（KF）难以直接应用，
因此在实际应用中，
通常采用扩展卡尔曼滤波（EKF）、
无迹卡尔曼滤波（UKF）等非线性滤波方法进行状态估计。
然而，这种非线性特性也使得直接方法在计算复杂度和稳定性方面
面临一定挑战。 \cite{HU2018755}

总体而言，
直接建模方法通过对导航参数进行直接估计，
能够较真实地反映系统动力学特性，
具有模型物理意义清晰的优点，
但由于系统非线性较强，
其在现代组合导航系统中的应用逐渐被误差状态建模的间接方法所取代，
后者在工程实践中表现出更好的数值稳定性和计算效率。

\subsubsection{间接建模方法}

下面对惯性导航系统的间接建模方法进行阐述。

在间接建模方法中，
系统并不直接对导航参数进行估计，
而是对导航误差进行建模与估计。
通过建立误差状态模型，
并利用滤波方法对误差进行估计，
再将估计误差反馈修正导航解，
可以有效提高系统稳定性与计算精度，
同时降低导航解的数值发散风险。

典型的惯性导航系统误差状态通常包括：
位置误差、速度误差、姿态误差以及惯性器件误差，
因此系统误差状态向量可表示为

\begin{equation}
\mathbf{x}
=
[
\delta \mathbf{r}^n,
\delta \mathbf{v}^n,
\boldsymbol{\phi},
\boldsymbol{\varepsilon}^b,
\nabla^b
]^T
\end{equation}

其中，
$\delta \mathbf{r}^n$、
$\delta \mathbf{v}^n$、
$\boldsymbol{\phi}$
分别表示位置误差、速度误差以及姿态误差，
$\boldsymbol{\varepsilon}^b$
与 $\nabla^b$
分别表示陀螺和加速度计零偏误差。

误差状态方程的建立，
通常基于惯性导航基本方程，
通过对导航方程进行一阶线性化，
并忽略高阶小量，
可得到误差传播模型。
在小角度假设条件下，
姿态误差通常采用等效旋转矢量表示，
并利用方向余弦矩阵的小扰动展开，
从而得到位置、速度及姿态误差之间的耦合关系。

根据惯性导航系统的位置更新方程，
可以推导得到位置误差传播方程；
同理，
由速度更新方程可得到速度误差传播方程；
而由姿态更新方程，
结合小角度误差模型，
可以得到姿态误差传播方程。
此外，
考虑惯性器件误差通常具有随机游走或一阶高斯–马尔可夫过程特性，
可建立陀螺及加速度计零偏误差模型。

综合上述各误差传播关系，
惯性导航系统误差状态模型
可统一表示为如下形式：

\begin{equation}
\begin{aligned}
\delta \dot{\mathbf{r}}^n
&= -\boldsymbol{\omega}_{en}^n \times \delta \mathbf{r}^n + \delta \mathbf{v}^n, \\
\delta \dot{\mathbf{v}}^n
&= \mathbf{C}_b^n \mathbf{f}^b \times \boldsymbol{\phi}
	- (2\boldsymbol{\omega}_{ie}^n + \boldsymbol{\omega}_{en}^n) \times \delta \mathbf{v}^n \\
&\quad - (2\delta \boldsymbol{\omega}_{ie}^n + \delta \boldsymbol{\omega}_{en}^n) \times \mathbf{v}^n
	+ \delta \mathbf{g}^n + \mathbf{C}_b^n \nabla^b, \\
\dot{\boldsymbol{\phi}}
&= -\boldsymbol{\omega}_{in}^n \times \boldsymbol{\phi}
	+ \delta \boldsymbol{\omega}_{in}^n - \mathbf{C}_b^n \boldsymbol{\varepsilon}^b, \\
\dot{\boldsymbol{\varepsilon}}^b
&= -\frac{1}{T_{\varepsilon}}\boldsymbol{\varepsilon}^b + \mathbf{w}_{\varepsilon}, \\
\dot{\nabla}^b
&= -\frac{1}{T_{\nabla}}\nabla^b + \mathbf{w}_{\nabla}.
\end{aligned}
\end{equation}

其中：

$\boldsymbol{\omega}_{ie}^n$
为地球自转角速度；

$\boldsymbol{\omega}_{en}^n$
为导航坐标系相对地球坐标系角速度；

$\boldsymbol{\omega}_{in}^n$
为导航坐标系相对惯性坐标系角速度；

$\mathbf{C}_b^n$
为体坐标系到导航坐标系的方向余弦矩阵；

$\mathbf{f}^b$
为加速度计测得的比力；

$\delta \mathbf{g}^n$
为重力误差项；

$\mathbf{w}_{\varepsilon}$ 与
$\mathbf{w}_{\nabla}$
分别表示陀螺与加速度计误差驱动噪声。

\subsection{GNSS的数学模型}

\subsection{误差的滤波方法}

KF
EKF
CKF
FGO方法
UKF
最大后验和极大似然估计

\section{本文研究意义}

\section{论文结构}

